\section{Compatibility Relation Examples}~\label{sec:comp-ex}

Consider these two Protobuf descriptors

\begin{figure}[H]
	\centering
	\begin{minipage}[bt]{0.43\textwidth}
		\begin{lstlisting}[language=proto]
// v1.proto -- A base proto for the
// example software. Tracks some
// basic information about an item,
// with the ledger itself basically
// being a list of items.
syntax = "proto3";

message Location {
  string loc = 2;
}

message Item {
  string name = 2;
  string description = 3;
  int32 value = 4;
  Location loc = 5;
}

message Ledger {
  int32 version = 1;
  string world_name = 2;
  repeated Item log = 3;
}\end{lstlisting}
	\end{minipage}
	\hspace{0.5cm}
	\begin{minipage}[bt]{0.48\textwidth}
		\begin{lstlisting}[language=proto]
// v2.proto -- build off v1.proto by
// adding some extra fields to the
// item, such as item level, and
// split descriptions and locations
syntax = "proto3";

message Location {
  string loc = 2;
}

message Item {
  string name = 2;
  string public_description = 3;
  string private_description = 7;
  int64 value = 4;
  int32 level = 6;
  Location origin_loc = 5;
  Location current_loc = 8;
}

message Ledger {
  int32 version = 1;
  string world_name = 2;
  repeated Item log = 3;
}\end{lstlisting}
	\end{minipage}

	\caption{Example Protobuf descriptors for a hypothetical piece of software
      which tracks TTRPG items, showing the original version on the left and an
      updated version on the right.}
	\label{fig:proto-ex}
\end{figure}

Clearly these descriptors are different, but are they \emph{compatible}?
Evidence of their compatibility is a derivation of the compatibility relation.
First, use the Message Relation (Section~\ref{sec:msg-rel}) to show that the base
descriptors are compatible. Since the \texttt{Location} and \texttt{Ledger}
messages are the same, one application of the reflexivity rule shows these to be
compatible. Focus on the \texttt{item} message then. 
\begin{align*}
  &\msg{\emptyset}{f} \\
  &\left\{f(3) = (\lb{description}, \str) \right\} \\
  &\msg{\emptyset}{f[3 \mapsto \field{public\_description}{\str}]} = \msg{\emptyset}{f_1} \\
  &\left\{ \ints \propto \intl \right\} \\
  &\msg{\emptyset}{f_1[4 \mapsto \field{value}{\intl}]} = \msg{\emptyset}{f_2} \\
  &\left\{ f_2(5) = \field{loc}{\mathtt{Location}} \right\} \\
  &\msg{\emptyset}{f_2[5 \mapsto \field{current\_loc}{\mathtt{Location}}]} = \msg{\emptyset}{f_3} \\
  &\left\{ 6 \not \in \dom{f_3} \right\} \\
  &\msg{\emptyset}{f_3[6 \mapsto \field{level}{\ints}]} = \msg{\emptyset}{f_4} \\
  &\left\{ 7 \not \in \dom{f_4} \right\} \\
  &\msg{\emptyset}{f_4[7 \mapsto \field{private\_description}{\str}]} = \msg{\emptyset}{f_5} \\
  &\left\{ 8 \not \in \dom{f_5} \right\} \\
  &\msg{\emptyset}{f_5[8 \mapsto \field{current\_loc}{\mathtt{Location}}]}
\end{align*}

If we had concrete values for these descriptors, it would also be possible to
show that those two values were related using the value relation.

% Local Variables:
% citar-bibliography: ("../pollux.bib")
% TeX-master: "../pollux.tex"
% jinx-local-words: "protobuf"
% End:
